\section{Basic closure properties}
\label{sec:ch34-closure}

\begin{warning}{Primary-source import; Rethlas gate}
The next theorem is an exact source-aligned import from
\cite[Definition~4.1, Theorems~4.5 and~4.7, Lemmas~4.6 and~4.9,
Corollary~4.17, Theorem~4.22, and Corollary~4.23]{BJ22}.  The present audit
found no exact-current Rethlas receipt for this repaired block.  Accordingly
none of its clauses is presented as newly machine-verified mathematics, and
no extension beyond the cited clauses may be inferred without a separate
source proof or an exact-current receipt.
\end{warning}

\begin{theorem}
\label{thm:bj-basic-closure}
Let \(\theta,\theta'\in N^1(X)\) be arbitrary numerical classes.  No
positivity hypothesis is needed: by \Cref{lem:ch34-psh-nonempty} the
assertions below are vacuous unless the relevant classes are
pseudoeffective, and nefness is required only where a \emph{bounded} member
is asserted to exist, which is nowhere in this theorem.
\begin{enumerate}
\item If \(\phi\in\PSH_X(\theta)\) and \(c\in\mathbf R\), then
      \(\phi+c\in\PSH_X(\theta)\).
\item If \(\phi,\psi\in\PSH_X(\theta)\), then
      \(\max\{\phi,\psi\}\in\PSH_X(\theta)\).
\item If \(\phi\in\PSH_X(\theta)\) and
      \(\psi\in\PSH_X(\theta')\), then
      \(\phi+\psi\in\PSH_X(\theta+\theta')\), and
      \(a\phi\in\PSH_X(a\theta)\) for every \(a>0\).  Thus strictly
      positive linear combinations add the corresponding classes.  An
      endpoint with coefficient zero means that the corresponding summand
      is omitted; it is never interpreted as \(0\cdot(-\infty)\).
\item If \((\psi_j)_{j\in J}\) is a decreasing net in \(\PSH_X(\theta)\) and
      \(\psi:=\lim_j\psi_j\) satisfies \(\psi\not\equiv-\infty\) on any
      irreducible component, then \(\psi\in\PSH_X(\theta)\).
\item \emph{(Varying classes.)}  More generally, if \(\psi_j\in\PSH_X(\theta_j)\)
      with \(\theta_j\to\theta\) in \(N^1(X)\), the net \((\psi_j)\) is
      decreasing, and \(\psi:=\lim_j\psi_j\) is generically finite, then
      \(\psi\in\PSH_X(\theta)\).  Consequently, for a fixed function
      \(\psi\) the set \(\{\theta:\psi\in\PSH_X(\theta)\}\) is closed.
\item If \(f:Y\to X\) is a morphism of projective varieties and
      \(\phi\in\PSH_X(\theta)\), then either
      \(f^*\phi\in\PSH_Y(f^*\theta)\), or \(f^*\phi\equiv-\infty\) on an
      irreducible component of \(Y\).  In particular, restriction to every
      irreducible component \(X_\alpha\) belongs to
      \(\PSH_{X_\alpha}(\theta|_{X_\alpha})\).
\item\label{itm:bj-closure:determination} Every \(\phi\in\PSH_X(\theta)\) is finite on
      \(X^{\lin}\), hence on \(X^{\divis}\).  Moreover, for every upper
      semicontinuous \(\eta:X^{\An}\to\mathbf R\cup\{-\infty\}\),
      \[
       \phi\leq\eta\ \text{on }X^{\divis}
       \quad\Longleftrightarrow\quad
       \phi\leq\eta\ \text{on }X^{\An}.
      \]
      In particular, psh functions are determined by their divisorial
      values.
\end{enumerate}
\end{theorem}

\begin{proof}
(1) If \((\phi_i)\) approximates \(\phi\) as in \Cref{def:bj-psh}(3), then
\((\phi_i+c)\) is a decreasing net with the same classes, and
\(\phi_i+c\in\cH^{\mathrm{gf}}_{\mathbf R}(L_i)\): in a presentation of
denominator \(m_i\), replacing every weight \(\lambda_j\) by
\(\lambda_j+m_i c\) in \eqref{eq:gf-fs-presentation} adds \(c\) to the
function.  No positivity of \(\theta\) is used.

(2) The point that the two approximating families must have the same class
at each stage is essential.  The simultaneous approximation lemma
\cite[Lemma~4.9]{BJ22} gives a net of \(\mathbf Q\)-line bundles \(L_i\) and
functions
\(\phi_i,\psi_i\in\cH^{\mathrm{gf}}_{\mathbf R}(L_i)\) such that
\(c_1(L_i)\to\theta\), \(\phi_i\searrow\phi\), and
\(\psi_i\searrow\psi\).  By
\Cref{lem:gf-fs-operations}\,(2),
\(\max\{\phi_i,\psi_i\}\in
\cH^{\mathrm{gf}}_{\mathbf R}(L_i)\), and these functions decrease
pointwise to \(\max\{\phi,\psi\}\).  The limit is generically finite
because \(\phi\) is, so it is \(\theta\)-psh.

(3) Choose approximating nets \((\phi_i)_{i\in I}\) and
\((\psi_{i'})_{i'\in I'}\).  On the directed product \(I\times I'\), the
net \(\phi_i+\psi_{i'}\) decreases to \(\phi+\psi\), and
\Cref{lem:gf-fs-operations}\,(1) places its terms in class
\(L_i+L'_{i'}\), whose Chern classes converge to \(\theta+\theta'\).
The limit is generically finite since both summands are finite at every
trivial valuation.  This proves the sum assertion.  Positive rational
scalars follow from
\Cref{prop:ch34-fs-operations}\,\cref{itm:ch34-ops:rational}.  For
\(a\in\mathbf R_{>0}\), subtract the finite number \(\sup\phi\) so that
\(\phi\leq0\), choose rational \(a_m\uparrow a\), and apply~(5) to the
decreasing sequence \(a_m\phi\in\PSH_X(a_m\theta)\); then restore the
finite constant \(a\sup\phi\) using~(1).  All scalar multiplications in
this argument have a strictly positive scalar.  Together with the sum
assertion this proves the stated positive-linear-combination rule.  This is
\cite[Theorem~4.7\,(i)--(ii)]{BJ22}.

(5) This is precisely property \emph{(PSH2)} of
\cite[Theorem~4.5]{BJ22}.  We spell out the diagonal issue that its proof
resolves.  Apply \cite[Lemma~4.6]{BJ22} with
\[
 \begin{aligned}
  K&=X^{\An},& S&=\Pic(X)_{\mathbf Q},
    &F_L&=\cH^{\mathrm{gf}}_{\mathbf R}(L),\\
  T&=N^1(X),& \rho(L)&=c_1(L).
 \end{aligned}
\]
Here \(K\) is compact, \(T\) is finite dimensional and hence metrizable,
and every \(F_L\) is stable under addition of a real constant by
\Cref{lem:gf-fs-operations}\,(3).  The lemma constructs the diagonal on
triples \((j,i,\varepsilon)\), ordered simultaneously by the outer index,
the pointwise order of the shifted approximants, and the class error.  That
preorder is directed by the compactness argument in the lemma and produces
one decreasing GF net whose classes converge to \(\theta\) and whose
pointwise limit is \(\psi\).  This is the required condition in
\Cref{def:bj-psh}(3).  In particular, merely ordering the disjoint union of
the inner index sets does \emph{not} prove directedness.  The closedness
assertion is the special case in which \(\psi_j=\psi\) for all \(j\).

(4) is the fixed-class special case of~(5).

(6) Write \(\phi_i\searrow\phi\) as in \Cref{def:bj-psh}.  Pullback of a
GF presentation is again GF unless all its sections vanish identically on
some component of the source, by \Cref{lem:gf-fs-operations}\,(4).  Thus,
if \(f^*\phi\) is generically finite, every \(f^*\phi_i\) is generically
finite, \(f^*\phi_i\searrow f^*\phi\), and
\(c_1(f^*L_i)\to f^*\theta\); hence
\(f^*\phi\in\PSH_Y(f^*\theta)\).  Otherwise the alternative in the
statement holds.  Restriction to \(X_\alpha\) is generically finite by
definition.  This is \cite[Theorem~4.7\,(iii)--(iv)]{BJ22}.

(7) Finiteness on \(X^{\lin}\) is
\cite[Corollary~4.17]{BJ22}.  The comparison statement is
\cite[Theorem~4.22]{BJ22}, and divisorial determination is its
\cite[Corollary~4.23]{BJ22}.
\end{proof}

\begin{warning}{Suprema are not yet available}
\Cref{thm:bj-basic-closure} contains no statement about \emph{increasing}
families, and this omission is the central open point of the trivially
valued theory as it stands at this stage of the book.  Given a nonempty
family \((\phi_\alpha)\subset\PSH_X(\theta)\) that is bounded above, the
regularized supremum \((\sup_\alpha\phi_\alpha)^*\) need not be visibly
\(\theta\)-psh: the definition supplies decreasing nets, and there is no
formal manoeuvre converting an increasing family into one.  That
\((\sup_\alpha\phi_\alpha)^*\) is nevertheless \(\theta\)-psh, for \(X\)
integral and unibranch, is the envelope theorem, stated and proved in
\Cref{chap:bj-envelope-statement}.  Until then, no argument in this book
uses closure under suprema, and the reader should treat every appeal to an
upper envelope before that chapter as a forward reference to be discharged
there.
\end{warning}
